% Folder 93, batch 10 — archivists' pages 181–200.
% Pass: Opus 5.5 (claude-opus-5-5), 2026-10-10 — first pass, unchecked against the pages by a human.
%
% Manuscript notes in Grothendieck's hand throughout. Page 187 is a cover sheet
% inscribed in his hand; it carries no \page{} and its words give the section
% heading. Page 186 was scanned upside down and is cancelled by long strokes.
% Page 199 carries only a short struck passage at its head.
\documentclass[11pt,a4paper]{article}
\input{../preamble/grothendieck}

\folder{93}
\batch{10}
\pages{181}{200}
\foldertitle{Connexions de Gauss-Manin et équations de Picard-Fuchs. Opération de Cartier : copies de tapuscrit annoté (s.d.), tiré à part (1981), notes manuscrites (s.d.).}
\dating{[1972]-1981}
\watermark{Édition de démonstration}

\begin{document}

\section{Opération de Cartier en dimension relative un}

\page{181}
\note{La page commence au milieu d'une phrase : l'argument vient d'avant le lot.}
… en question \ill{} sur \emph{Lie} et l'opération des puissances
\uncertain{$p$-ièmes} (c'est-à-dire dans le livre de
\uncertain{Demazure-Gabriel} ?).
\marginal{?}
L'isom transposé est \uncertain{donné} par l'opération de Cartier sur les
formes différentielles.

\emph{⑦ Cas \add{lisse} de dim relative un : un accouplement remarquable.}

\struck{Soit} Plus généralement, en \uncertain{dim relative} $1$,
\struck{$Q = f_{X/S*}(\mathcal{O}_X)/\mathcal{O}_{X^{p/S}}$}
supposons que $X$ sur $X^{p/S}$ \struck{\ill{}}, que $\Omega^1_{X/S}$
($=\Omega^1_{X/X^{p/S}}$) soit \emph{lisse de rang 1}, ce qui signifie aussi que
$X$ a loc. une $p$-base à $1$ \uncertain{élément} sur $S$. Posons
\[
  Q = f_{X/S*}(\mathcal{O}_X)/\operatorname{Im}\mathcal{O}_{X^{p/S}}
    = f_{X/S*}(\mathcal{O}_X)/\Lambda
    \simeq \mathcal{O}_X/f^{-1}\mathcal{O}_S[\mathcal{O}_X^{p}]
\]
qui est un module loc. libre de rang $p-1$ sur $S$. On va définir un
accouplement \add{$\Lambda$-bilinéaire} canonique
\[
  (*)\qquad Q\times Q \longrightarrow
  \Omega^1_{X^{p/S}/S}\otimes_{\mathcal{O}_{X^{p/S}}}\Lambda
\]
($=\Omega^1_{X^{p/S}/S}$ si p.ex. $X/S$ lisse), ou plutôt, ce qui revient au
même, une forme $\Lambda$-bil.
\[
  \varphi : f_{X/S*}(\mathcal{O}_X)\times f_{X/S*}(\mathcal{O}_X)
  \longrightarrow \Omega^1_{X^{p/S}/S}\otimes_{\mathcal{O}_{X^{p/S}}}\Lambda
\]
qui est nulle sur $\Lambda\times f_{X/S*}(\mathcal{O}_X)$ et sur
$f_{X/S*}(\mathcal{O}_X)\times\Lambda$.

\page{182}
On prend
\[
  \varphi(f,g) = \widetilde{C}(f\,dg)
\]
a) $\Lambda$-bilinéaire : clair.

b) nulle si \struck{$f$ scalaire de $\Lambda$} $f=1$ ou $g=1$ : clair, car
$\widetilde{C}$ est nulle sur les diff. exactes.

c) $\varphi$ est antisymétrique, i.e.
$\left\{\begin{array}{l}\text{alterné si } p\neq 2\\ \text{symétrique si } p=2\end{array}\right.$
i.e. $\varphi(f,g)+\varphi(g,f)=0$ i.e.
$\widetilde{C}(f\,dg+g\,df)=0$, en effet $f\,dg+g\,df = d(fg) = 0$.
\note{Ainsi sur la page ; c'est $\widetilde{C}(d(fg))$ qui est nul.}

\emph{Proposition.} L'accouplement $(*)$ est une dualité parfaite de
$\Lambda$-modules.

On peut supposer qu'il existe une $p$-base relative $\{x\}$, et on a une base
de $Q$ sur $\Lambda$ formée des \struck{$x, x^2, \dots, x^{p-1}$} $x^i$,
$1\leq i\leq p-1$. Or on a
\[
  \widetilde{C}(x^i\,dx^j) = \widetilde{C}(j\,x^{i+j-1}\,dx) =
  \begin{cases} 0 & \text{si } i+j\neq p-1\\ j & \text{si } i+j = p-1\end{cases}
\]
\note{Ainsi sur la page ; le calcul donne la condition $i+j-1=p-1$, soit $i+j=p$, ce que confirme la matrice antidiagonale qui suit.}
donc la matrice de la forme $\varphi$ est
\[
  \begin{pmatrix} & & p-1\\ & \cdots & \\ 1 & & \end{pmatrix}
  \qquad\text{(antidiagonale $1, 2, \dots, p-1$, zéros ailleurs),}
\]
son déterminant, au signe près, est donc $1.2\dots(p-1)\in\mathbb{F}_p^*$, donc
inversible, O.K.
\marginal{\uncertain{$\Delta = -1 \in \mathbb{F}_p^*$}}

\emph{Corollaire.} \struck{$\det Q$} On a un isom canonique
\[
  Q \xrightarrow{\ \sim\ } \check{Q}\otimes\bigl(\Omega^1_{X^{p/S}/S}\otimes\Lambda\bigr)
\]
d'où un isom canonique
\[
  \det Q \simeq \det\check{Q}\otimes
  \bigl(\Omega^{1\ \otimes(p-1)}_{X^{p/S}/S}\otimes\Lambda\bigr)
\]

\page{183}
et par suite un isom canonique
\[
  (**)\qquad (\det Q)^{\otimes 2} \simeq
  \Omega^{1\ \otimes(p-1)}_{X^{p/S}/S}\otimes\Lambda
\]

\emph{Cas 1} $p=2$, on trouve
\[
  \text{\struck{$\det Q^{\otimes}\simeq$}}\ \Omega^1_{X^{p/S}/S}\otimes\Lambda
  \simeq (\det Q)^{\otimes 2}
\]
donc $\Omega^1_{X^{p/S}/S}\otimes\Lambda$ \add{et $X$ réduit géom.,
$\Lambda\simeq\mathcal{O}_{X^{p/S}}$} est un \emph{carré}, (ceci implique que
$\Omega^1_{X/S}$ est lui-même un carré, \ill{} si $S=\operatorname{Spec}(k)$,
$k$ \ill{}) \struck{\ill{}} \add{\ill{} si $X$ lisse sur $k$} communique un
\uncertain{carré}. \struck{\ill{}} \ill{} \uncertain{comme} on \uncertain{voit}
\struck{\ill{}} par l'exemple de \uncertain{Serre-Tate-Fröhlich} : il suffit de
prendre l'exemple d'une courbe \add{lisse propre} rationnelle sans point
rationnel \ill{} --
\marginal{?}
\struck{\ill{}} il y a un \ill{} sur \uncertain{corps imparfait} du cas 2.
\note{Passage serré, surchargé de ratures et d'ajouts interlinéaires ; seules les lectures indiquées sont proposées.}

\emph{Cas 2} $p\neq 2$, donc $\frac{p-1}{2}$ est un entier. On aimerait
\uncertain{trouver} un isom can.
\[
  \psi : \det(Q) \simeq
  \Omega^{1\ \otimes(\frac{p-1}{2})}_{X^{p/S}/S}\otimes\Lambda
  \qquad\Bigl[\simeq \Omega^{1\ \otimes(\frac{p-1}{2})}_{X/S}
  \otimes_{\mathcal{O}_X}\Lambda\Bigr].
\]
\uncertain{qui serait} un \uncertain{relèvement} de $(**)$ \ill{} de
$\mu_2\simeq\mathbb{Z}/2\mathbb{Z}$ [\uncertain{Sinon}, il y aurait un
\uncertain{revêtement} \ill{} de $X^{p/S}$, canoniquement associé à la
situation]. Regardons la base de $Q$, $(\dot{x}^i)_{1\leq i\leq p-1}$,
\struck{et comme} d'où une base de $\det Q$,
$\dot{x}\wedge\dot{x}^2\wedge\dots\wedge\dot{x}^{p-1}$, et la valeur de
l'isom $(**)$, est
$\Delta\bigl(d\tilde{x}^{\otimes(p-1)}\otimes\Lambda\bigr)$, où $\Delta$ est le
déterminant de la matrice \uncertain{écrite} plus haut, savoir
\[
  (-1)^{(p-1)-1}\ 1.2\dots(p-1) = \text{\struck{$-(1.2\dots(p-1))$}}\,.
\]
Comme $\mathbb{F}_p^*$ est un groupe cyclique d'ordre \struck{pair} $p-1$, le
produit des \uncertain{éléments} est $-1$, donc on trouve $\Delta = 1$,

\page{184}
Ceci suggère \struck{d'associer à $\dot{x}$ … $\dot{x}$} de définir $\psi$ par
\[
  (***)\qquad \boxed{\ \Psi(\dot{x}\wedge\dot{x}^2\wedge\dots\wedge\dot{x}^{p-1})
  = \widetilde{dx}^{\,\otimes\frac{p-1}{2}}\otimes\Lambda\ }
\]
Ce qui demande une vérification \uncertain{courte} de l'isom $(**)$, et
\uncertain{d'autre part} que $\psi$ est indépendant du choix de la $p$-base
$\{x\}$. \struck{\ill{} une autre $p$-base,}
\struck{$y = a_{p-1}x^{p-1}+\dots+a_0$ $(a_i\in\Lambda)$, d'où}
\struck{$dy = [(p-1)a_{p-1}x^{p-2}+\dots+a_1]\,dx$, donc on a}
\struck{$(p-1)a_{p-1}x^{p-2}+\dots+a_1\in B^*$}
\note{Ce calcul, depuis « une autre $p$-base », est barré par de longs traits obliques.}
Or, \ill{} de l'\ill{} : l'autre \ill{} \uncertain{façon} \ill{}
\struck{\uncertain{complexe}} \add{\uncertain{simplicial}} des schémas
$P$ [\ill{} $Z$ \ill{} $(X^{p}, \Lambda)$] \ill{} des $p$-bases
relatives, dans le revêtement trivial à deux \uncertain{feuillets}
$(\mathbb{Z}/2\mathbb{Z})_Z$ de $Z$. \struck{\ill{}} Or \struck{\ill{}}
$P \to Z$ \ill{} : fibres connexes, et on voit donc que
\struck{cette application est} \ill{} une \uncertain{application} est
\uncertain{constante} \ill{} fibres : fibres, ce qui \uncertain{donne} ce qu'on
voulait …

\uncertain{Application} : la \uncertain{définition} cohomologique globale
(\uncertain{Serre} \ill{} \uncertain{Sém.} vol.~1)

\marginal{?}
\uncertain{Généralisation} des relations précédentes à la dim supérieure ?

\page{185}
Cas $X\xrightarrow{f}S$ propre \add{lisse} de dim relative $1$. On a un
diagramme autodual [dualité = symétrie p.r. à un \uncertain{centre}]

\begin{tikzcd}[column sep=small, row sep=small, nodes={font=\scriptsize}]
 & & 0 \arrow[d] & & \\
 & & R^1f_*(\mathcal{O}_{X/S}) & & \\
0 \arrow[r] & R^1f_*(\mathcal{O}_X)^{(p)} \arrow[r, "\alpha"] \arrow[ur, "F\ (\mathrm{Frob})"] & \mathcal{H}^1_{DR} \arrow[r, "\beta"] \arrow[u] & f_*(\Omega^1_{X/S})^{(p)} \arrow[r] & 0\\
 & & f_*(\Omega^1_{X/S}) \arrow[u] \arrow[ur, "C\ (\text{op. de Cartier})"'] & & \\
 & & 0 \arrow[u] & &
\end{tikzcd}

\note{Disposition reproduite d'après la page : la colonne monte de $0$ à $0$ ; la flèche du haut, de $\mathcal{H}^1_{DR}$ vers $R^1f_*(\mathcal{O}_{X/S})$, et la flèche de $0$ vers celui-ci sont des traits sans tête nette sur le manuscrit.}

\marginal{\emph{ligne et colonne exactes}}

$*$) $\mathcal{H}^1_{DR} = \mathcal{H}^1_{DR}(X/S) = R^1f_*(\Omega^{\bullet}_{X/S})$

$**$) colonne : est la filtration de $\mathcal{H}^1$ par $\mathcal{H}^{1,0}$ et
$\mathcal{H}^{0,1}$

$***$) $\alpha$ déduit par passage au quotient de
$F_{\mathcal{H}^1} : \mathcal{H}^{1(p)}\to\mathcal{H}^1$

$****$) \uncertain{$\beta = {}^t\alpha$}, déduit de
\uncertain{${}^tF_{\mathcal{H}^1}$} par restriction

$*****$) $F$ et $C$ sont \struck{duaux} \add{transposés} l'un de l'autre
\uncertain{via} la dualité de Poincaré et Cartier.

\note{L'encadré qui suit est barré par trois longs traits obliques.}
\struck{On a $\beta\alpha = 0$, et on a équivalence entre les conditions}

\struck{(i) $\alpha$ injectif}

\struck{(ii) $\beta$ surjectif}

\struck{(iii) la suite}

\struck{$0\to R^1f_*(\mathcal{O}_X)^{(p)}\xrightarrow{\alpha}\mathcal{H}^1_{DR}$}
\struck{$\xrightarrow{\beta} f_*(\Omega^1_{X/S})^{(p)}\to 0$ est exacte}

\struck{(iii bis) $\operatorname{Ker}\beta = \operatorname{Im}\alpha$, i.e.}
\struck{$\operatorname{Ker}\beta = (\operatorname{Ker}\beta)^{\perp}$ (ou $\supset$)}
\struck{i.e. $\operatorname{Im}\alpha = (\operatorname{Im}\alpha)^{\perp}$ (ou $\subset$)}

\struck{(iii ter) $\dim\operatorname{Ker}\beta = \dim\operatorname{Im}\alpha$}

\struck{(iv) $\operatorname{rang}\alpha = g$, i.e. $\operatorname{rang}\beta = g$,}
\struck{i.e. $\operatorname{rang}F_{\mathcal{H}^1} = g$}

\struck{Pour ceci, il suffit que}

\emph{NB} on a \uncertain{toujours}
$\left\{\begin{array}{l}\operatorname{Ker}\alpha = (\operatorname{Im}\beta)^{\perp}\\
\operatorname{Im}\alpha = (\operatorname{Ker}\beta)^{\perp}\end{array}\right.$
\struck{\ill{}}

\page{186}
\note{Page entière annulée par trois longs traits courbes.}

\begin{tikzcd}
X \arrow[r, "f^{(i)}"] & X^{[i]}\\
T_0 \arrow[u, "\varphi_0"] \arrow[r, "i"'] & T \arrow[u, "u"']
\end{tikzcd}

tels que $\left\{\begin{array}{l}\varphi_0\in\mathcal{S}_0\\ i\in\mathcal{N}\end{array}\right.$

Considérons \struck{\ill{}} plus généralement une \uncertain{famille
d'immersions} $\Gamma : X\to Z$ \add{de $(\mathrm{Sch})/S$}, et la catégorie
analogue formée des

\begin{tikzcd}
X \arrow[r, "\delta"] & Z\\
T_0 \arrow[u, "\varphi_0"] \arrow[r, "i"'] & T \arrow[u, "u"']
\end{tikzcd}

\ill{} \ill{} \ill{} \uncertain{désignera} $C/\delta$. Si $Z_{(n)}$ est le
$n$-ième \uncertain{voisinage} infinitésimal de $X$ dans $Z$, \struck{\ill{}}
et $\delta_n : X\to Z_{(n)}$, alors $C/\delta$ est la \emph{limite} des
sous-catégories $C/\delta_n$, qui \ill{}
\struck{\ill{}} [\uncertain{si} \ill{} des \uncertain{voisins}
\struck{$X$} $\delta_n$, $Z_n\in\mathcal{B}S$]
\struck{et que $X\hookrightarrow Z_n$ \ill{} loc.}
\struck{\ill{} $X$ \ill{} $(C)$ une restriction sur $X$} ],

\section{Connexions et stratifications intégrables : (suite, problèmes divers)}
\note{Titre tiré de la feuille de garde p. 187, de sa main ; la feuille ne porte rien d'autre.}

\page{188}
$X/S$ schéma sur $S$.

\marginal{Il serait plus naturel à certains égards de travailler avec les PD-connexions [c'est pareil dans cas lisse]}
1) Un module \add{qu. coh.} $E$ sur $X$, avec connexion rel/$S$, est dit
\emph{trivial} si $E\simeq f^*(F)$, $F$ un module \add{qu. coh.} sur $S$
[isom respectant connexions]. \emph{NB} Si $X\to S$ est
\struck{\uncertain{connexe}} \add{lisse surjectif} \uncertain{fpqc} et à fibres
géométriquement connexes, \add{$S$ de car.~$0$, \uncertain{donc}} le foncteur
$F\mapsto f^*(F)$, des modules qu. coh. sur $S$ dans les modules qu. coh. sur
$X$ à connexion rel/$S$, est pleinement fidèle [\uncertain{même} \uncertain{aussi}
si $F$ de présentation finie].

On dit que $E$ est loc. \emph{trivial} au \emph{Zar.-intégrable} si $E$ est
trivial localement sur $(X\to S)$ pour la topologie de Zariski. On dit que
$E$ est \emph{intégrable}, si la même condition est satisfaite, loc. pour top.
ét.

2) Supposons $S$ de car. $p>0$, et que $X$ sur $S$ admette loc. une $p$-base.
Alors \struck{$E$ est loc. trivial} conditions suivantes sur $E$ qu. coh. sont
équivalentes :

a) $E$ est \emph{Zar. intégrable}

b) $E$ est intégrable

c) $E$ est à courbure et $p$-courbure nulle.

d) $E$ est \emph{trivial} relativement à
$X\to\bigl(X^{p/S}, \mathcal{O}_{X^{p/S}}/J\simeq\Lambda\bigr)$

De plus, le foncteur $F\mapsto f'^{*}_{X/S}(F)$ des Modules
\uncertain{qu. coh.} sur
$X'^{/S'} = \bigl(X^{p/S}, \mathcal{O}_{X^{p/S}}/J = \Lambda\bigr)$ dans les
modules \add{qu. coh. sur $X$} à connexion intégrable rel/$S$, est une
équivalence de catégories.

\page{189}
3) Soit $G$ un groupe \struck{fini} étale sur $S$, et considérons un
$G_X$-torseur $P$ sur $X$ (à droite). \struck{Pour tout} \add{Considérons une}
représentation de $G$ dans un module qu. coh. $F$ sur $S$. Alors $G_X$ opère
sur $F_X$, et on constate qu'il opère en laissant invariante la connexion
canonique de $F_X$. Comme les Modules à connexion intégrable \add{rel/$S$} sur
des \uncertain{étales} de $X$ forment un champ \uncertain{(considéré comme
\add{fibré} au-dessus de $S$)}, on voit que $G_X$ opère sur \uncertain{la}
section \add{$(F_X)$} du dit champ, qui se \uncertain{peut} tordre par le
$G_X$-torseur $P$. On trouve donc un module à connexion intégrable
$P_X\times^{G_X}F_X$.

\marginal{En fait, \uncertain{si} \ill{} \uncertain{un argument}, \ill{} une \emph{stratification} \emph{intégrable} sur $P_X\times^{G_X}F_X$. Cela \uncertain{revient} au \ill{} si $X/S$ lisse et $S$ de car.~$0$.}
\marginal{\uncertain{Donnons} d'abord \ill{} les \uncertain{\ill{}-modules}, \uncertain{gerbes} \ill{}}

Cette façon de construire des \struck{connexions} \add{stratifications}
intégrables mérite \uncertain{une} \struck{\ill{}} \uncertain{étude} \ill{}
\uncertain{universelle} dans une certaine mesure.
\struck{$X\to S$ soit plat, \ill{} \add{\uncertain{lorsque} \ill{}} à fibres
\ill{}, et qu'il ait une section $e$,}
\struck{Supposons} \add{\uncertain{géométriquement} \ill{}}
\note{Encadré raturé de traits obliques ; on n'en donne que les fragments lisibles.}

Supposons $S$ \uncertain{réduit} au spectre d'un corps $k$, et que $X$ est
géom. connexe et loc. de présentation finie sur $k$, muni d'un pt $x\in X$
\struck{\ill{}} \add{rationnel sur} $k$, on peut former le
\uncertain{pro-groupe} \add{(\uncertain{strict})} (lorsque $X$ n'est pas
normal) \uncertain{proprement dit} $\underline{\Pi}_1(X,x)$, et la catégorie des
Modules à \uncertain{stratification} intégrable sur $X$ doit être
équivalente à celle des représentations linéaires de dim finie de
$\underline{\Pi}_1(X,x)$. On devrait pouvoir généraliser au cas d'une base
quelconque $S$, en se

\page{190}
donnant un \uncertain{progroupe} $\mathbb{L}$ de nbs premiers, et supposons
que l'on a un \struck{\ill{}} \add{\uncertain{pro-groupe}}
$\underline{\Pi}_1^{\mathbb{L}}(X,e)$, obtenu à l'aide d'une section $e$ de
$X/S$ [plus généralement, une \uncertain{multisection} $E$ de $X$ finie, et
radiciel surjectif sur $S$] (on suppose $X\to S$ plat de présentation finie,
\add{à fibres géom. connexes,} $S$ loc. noeth. (?)) : la \add{cat. des}
représentations linéaires de $\underline{\Pi}_1^{\mathbb{L}}(X,e)$ dans des
modules cohérents sur $S$ doit être équivalente à celle des
\struck{représentations} Modules coh. sur $X$ à \struck{strat.} \add{stratif.}
intégr. rel/$S$, qui sur chaque fibre \add{géométrique} correspondent à une
représentation du groupe fondamental qui \struck{soit} se factorise via un
$\mathbb{L}$-groupe. Plus généralement, si on a un \struck{qu'on se pose}
pro-(groupe \struck{\ill{}} \add{étale}) \struck{\uncertain{loc. constant}}
strict $\Pi$ sur $S$, et un pro-torseur $P$ sous $\Pi_X$ sur $X$, tel que
\add{sur chaque} fibre géom. \struck{fibre} $\Pi_{\bar{s}}$ apparaisse comme
\emph{quotient} du groupe fondamental \uncertain{étale} de $X_{\bar{s}}$, alors
on \struck{\uncertain{doit}} trouve une équivalence entre la catégorie des
rep. (continues) de $\Pi$ dans des modules coh. sur $S$, et celle des Modules
coh. à stratification intégrable sur $X$ qui, sur chaque fibre géom., restent
bien être \uncertain{\ill{}} [\uncertain{définis}] à $\Pi_{\bar{s}}$
(regardé comme quotient du groupe fondamental de $X_{\bar{s}}$) …
\marginal{?}

\page{191}
\marginal{loc. de t.f. suffit ?}
4) Soit $X$ de type fini sur un corps, $X$ \struck{\ill{}} \emph{géom.
unibranche}, on veut étudier les modules inversibles sur $X$ à stratification
intégrable, \struck{\ill{}}. Soit $L$ un module \add{inversible} stratifié sur
$X$. Je dis que la stratification est intégrable ssi $\exists$ entier $n>0$,
\add{(\uncertain{premier à la car.})} tel que $L^{\otimes n}$ soit trivial
comme module stratifié [Le « si » est valable sans hyp. géom. unibranche],
i.e. ssi \struck{$L$ \ill{}} $L$ donne un élément d'ordre fini \add{premier à
$p$} dans le groupe des modules inversibles stratifiés. À fortiori, une
\add{premier à $p$} condition nécessaire pour qu'il existe un entier $n>0$
\add{premier à $p$} tel que $L^{\otimes n}$ admette une section partout $\neq
0$, i.e. que $L$ donne un élément d'ordre fini \add{premier à $p$} dans
$\operatorname{Pic}(X)$. Pour exprimer cette condition \struck{\ill{}}
\uncertain{maintenant}, il reste à exprimer que $L^{\otimes n}$ est à
\struck{stratification} \add{stratification} intégrable. On sait que les
\add{stratifications} \struck{\ill{}} de rang \uncertain{un} correspondent,
pour les \struck{connexions} \add{celles qui sont} intégrables. Or, si $X/k$
\ill{} sur $\mathcal{O}_X$, \struck{\ill{}} \add{\uncertain{de car. nulle}}
\marginal{ici de car. $0$ !}
correspondent aux stratifications sur $\mathcal{O}_X$ \uncertain{i.e.} aux
connexions à courbure nulle, i.e. aux formes \add{différentielles} fermées
\uncertain{sur} $X$. Dire que la stratification est triviale
\struck{\ill{}, la \ill{}} revient à dire que \struck{c'est} \emph{une
différentielle logarithmique}. Dire qu'elle est intégrable revient donc à dire
qu'il y a un entier $n>0$ tel que $n\omega$ soit une différentielle
logarithmique. Comme le groupe des différentielles logarithmiques est un
groupe isomorphe à un $L=\mathbb{Z}^r$ [\uncertain{d'après} Rosenlicht], il
est clair

\page{192}
que lorsqu'il n'est pas nul il n'est pas \emph{divisible}, \struck{\ill{}} au
groupe des formes fermées, multiples des logarithmiques, i.e. il y a des
\struck{\ill{}} connexions intégrables sur $\mathcal{O}_X$ qui ne sont pas
isomorphes à la connexion triviale, [en fait le groupe des classes
d'isomorphie de telles connexions intégrables est isomorphe à
$L\otimes\mathbb{Q}/\mathbb{Z}\simeq(\mathbb{Q}/\mathbb{Z})^r$]. Par contre,
si $X$ est \add{tel que} $\Gamma(X,\mathcal{O}_X^*)$ \struck{\ill{}} \add{forme
des sections horizontales)} complète, \struck{\ill{}} \add{une} connexion
intégrable sur $\mathcal{O}_X$ est triviale, \struck{la} une $1$-forme
différentielle qui est loc. ét. une différentielle logarithmique, est nulle.
\struck{La classification des} \add{\uncertain{\ill{}} à courbure nulle, i.e.}
\struck{Modules inversibles avec connexion} \add{stratification,} \ill{}
\marginal{\uncertain{différentielle} \ill{} \uncertain{sur} \ill{} \ill{} \uncertain{quelle} \ill{} \ill{} \uncertain{localement triviale} \ill{}}
\note{Passage encadré, en partie barré ; lecture fragmentaire.}

Revenons au cas où $k$ est de car. quelconque. Je dis \struck{\ill{}} que si
$\Gamma(X,\mathcal{O}_X)^*$ est formé de sections horizontales,
\struck{\ill{}} \uncertain{p.ex.} $X$ propre sur $k$ et séparable sur les
$\Lambda$\ill{}, il n'y a aucune autre \struck{\ill{}} stratification sur
$\mathcal{O}_X$ \add{intégrable} que la structure habituelle. En effet,
\add{(\uncertain{\ill{} de car. $p$})} si $\mathcal{O}_X^{\otimes n} =
\mathcal{O}_X$ a une section horizontale $\sigma$ \add{\uncertain{partout}
$\neq 0$} \struck{\ill{}}, \add{\uncertain{pour la}} stratification
$T^{\otimes n}$, \struck{qui est aussi horizontale}. \uncertain{Donc} la section
$\sigma_0 = 1$ de $\mathcal{O}_X$ satisfait $\sigma_0^{\otimes n} = \lambda\sigma$,
$\lambda\in\Gamma(X,\mathcal{O}_X)^*$, \add{\uncertain{\ill{}}} et
\uncertain{les} \add{\ill{}} horizontales, donc $\lambda\sigma$ aussi, donc
$\sigma_0^{\otimes n}$ horizontale, donc $\sigma_0$ aussi. Par suite, on voit
que la classification des \struck{\ill{}} Modules inversibles à connexion
\uncertain{intégrable} \emph{\uncertain{identique}} à celle des
\uncertain{\ill{}} \ill{}, dans ce cas, aux Modules inversibles \struck{\ill{}}
$L$ tels que $\exists\,n>0$, $n$ premier à $p$, avec $L^{\otimes n}\simeq
\mathcal{O}_X$ \ill{}
\note{La fin de la page déborde sur la suivante : la phrase continue p. 193.}

\page{193}
\emph{Questions}

$S$ schéma de type fini sur $\mathbb{Z}$, dominant $\mathbb{Z}$, $S$ intègre
pt géné. $\eta$

$X$ schéma lisse de type fini sur $S$.

$E$ module localement libre sur $X$, avec connexion à courbure nulle
relativement à $S$.

La \add{\uncertain{connexion}} \add{de $E_\eta$ sur $X_\eta$,} est dite
\emph{de niveau zéro} \add{ou \emph{intégrable}} \struck{\ill{}}, \struck{sur
$X_\eta$, \ill{}} \uncertain{lorsqu'il} \uncertain{peut} trouver une variété
finie étale surjective de $X_\eta$ sur laquelle $E_\eta$ devient localement
triviale (\emph{i.e.} « intégrable » au sens (Zar).). \struck{Une condition}
\uncertain{Je veux} \struck{\ill{}} \uncertain{on a \ill{} dire que sur une
extension} \struck{\ill{}}
\marginal{\uncertain{définition} à \ill{} \uncertain{\ill{}} \uncertain{qu'elle} \struck{\ill{}} \uncertain{il faut \ill{} localement} \struck{\uncertain{trivial}} \uncertain{(top. étale)} \ill{}}
dans $K$ de $k(\eta)$, $E_K$ sur $X_K$ soit associé à des repr. linéaires des
groupes fondamentaux des composantes connexes de $E_K$. Et si on
\uncertain{désire} $K=\mathbb{C}$, que les représentations des
$\pi_1$ \add{\uncertain{transcendants}} \struck{des composantes} \struck{\ill{}}
connexes, définies par $E_{\mathbb{C}}$, \struck{\ill{}} se \uncertain{fassent}
à travers des groupes finis.

Cette notion peut se varier aussi pour \uncertain{n'importe} quel schéma en
groupes de type f. $G$ sur $S$, en prenant sur $X$ un torseur $P$ sous $G_X$,
avec connexion à courbure nulle rel/$S$. La question qui suit \uncertain{peut}
se poser aussi pour de tels groupes.

Remarquons que si \struck{la} connexion de $E$ est intégrable, \uncertain{alors}

\page{194}
il existe un \struck{\uncertain{ouvert}} \add{voisinage} \struck{$U\neq\emptyset$
de $S$} de $\eta$ dans $S$, tel que pour tout $s\in U$, \add{\uncertain{hors de
car. $p>0$},} \struck{$E_s$ sur $X_s$, avec la connexion relative à
$k(s)$} \uncertain{la} $p$-courbure de $E_s$ sur $X_s/k(s)$ soit nulle.
\marginal{ou connexion intégrable}
[Cette dernière condition est d'ailleurs équivalente à la condition analogue,
\struck{\ill{}} où on se limite aux pts \emph{fermés} de $U$, mais peu
importe].

\emph{Questions} La \emph{réciproque est-elle vraie} ?

\emph{Cas \struck{\ill{}} où $E$ est de rang 1.} Notons que \struck{\ill{}}
$E_\eta$ intégrable signifie que $E$ est associé à une \uncertain{variété}
\struck{\ill{}} des $X_\eta$, i.e. \struck{\ill{} un module} $L$ ),
\uncertain{connexion} est munie d'une trivialisation de $E^{\otimes n}$, pour
$n\geq 1$ convenable. La question précédente se décompose en deux :

a) Si la connexion des $E_s$ \add{\uncertain{($s\in S$, de car.~$>0$)}} est
intégrable, est-il vrai que l'image de $E_\eta$ dans $\operatorname{Pic}(X_\eta)$
est de torsion ?

b) Si la réponse est affirmative, on est ramené dans le cas où $E\simeq
\mathcal{O}_X$ (isom ne respectant pas la connexion, à priori)

b) Si $\omega$ est une $1$-forme différentielle \add{[fermée]} sur $X/S$, et
si les $\omega_s$ sur les $X_s/s$ ($s\in S$, $s$ de car.~$>0$) sont loc. des
différentielles logarithmiques, alors est-il vrai que \add{(un multiple entier
de)} $\omega_\eta$ est une diff. logarithmique ?
\marginal{ce qui signifie par l'opération de Cartier $\widetilde{C}\omega_s = \widetilde{\omega}_s$}
\note{Deux paragraphes portent la lettre b) ; le premier est corrigé à la main sans qu'on puisse dire en quoi.}

\page{195}
(\uncertain{resp.} de rang unipotent nul, resp. affine, resp. unipotent).

\emph{NB} Les cas « \add{\uncertain{connexe}} affine », « unipotent » sont
\uncertain{triviaux}, et \uncertain{\ill{}} \uncertain{triviaux}
\struck{(la dimension, base au moins, \ill{} \ill{} \ill{}, plus gén.)}
(cas les \uncertain{plus} \ill{} dans \uncertain{\ill{}} \uncertain{autre}),
grâce au th. de Chevalley. Notons que dans le cas unipotent, l'hypothèse
signifie que $\omega(s)$ est dans $p^{n}$ ($p$ = car. de $s$). Le cas
« \uncertain{rang} unipotent nul » peut se reformuler ainsi, lorsque $E$ est
commutatif : supposons qu'il existe un entier $N$ tel que, pour tout pt
\emph{fermé} $s$ de $S$, $\omega(s)$ soit \struck{\ill{}} dans
$(p^N E)(s) = (p^N E_s)(s) = p^N E(s)$, i.e. qu'il soit divisible par $p^N$.
Alors $\omega(\eta)$ appartient à un \uncertain{sous-groupe}
\uncertain{unipotent} nul, donc \uncertain{algébrique} de $E_\eta$ de rang
\struck{\ill{}} unipotent nul, donc \uncertain{extension} d'un $\sigma$-groupe
\uncertain{abélien} \ill{} par un tore.
\marginal{dans ce cas, on voit \uncertain{que} $E_\eta$ \uncertain{\ill{}} \uncertain{de rang unipotent nul} \struck{\ill{}} plus \ill{} un \uncertain{extension} d'une \uncertain{variété} \ill{} V.A. par un \uncertain{tore}, et que \ill{} \uncertain{l'existence} \ill{} \uncertain{il faut} \ill{} \uncertain{telle que} \ill{} \uncertain{d'induction}}

Si ce résultat était acquis, \struck{\ill{}} les « \uncertain{cas} \ill{}
groupe » se ramènent au cas où $E$ est déjà une extension d'une schéma abélien
$A$ par un tore $T$.
\struck{\uncertain{L'autre} cas extrême \ill{} \ill{} de $E$ \ill{} \ill{}
p.ex. $E$ est \ill{} \ill{} d'une variété abélienne $A$ par un tore,
\ill{}, \ill{}, alors il faudrait \ill{} que si $\omega$ \ill{} tel que
$\omega(s)$}
\note{Passage encadré et barré de traits obliques.}
\uncertain{\ill{}} \uncertain{entraîne} aux trois autres cas, ce cas-ci
semble également intimement lié au \uncertain{cas} \uncertain{supposé}
$S$ de car. $p>0$ \uncertain{\ill{}} !

\page{196}
\struck{Considérons en particulier le cas où $X\to S$ est propre. Dans ce cas
$\underline{\omega} = f_*(\Omega^1_{X/S})$ faisceau cohérent sur $S$ et
\ill{} \uncertain{\ill{}} \uncertain{\ill{}}, on peut supposer que}
\struck{$\underline{\omega} = f_*(Z^1_{X/S})$ $= \operatorname{Ker}\bigl(f_*(\Omega^1_{X/S})$ $\to f_*(\Omega^2_{X/S})\bigr)$.}
\note{Encadré barré de traits obliques.}

Dans le cas où $X\to S$ est un schéma abélien, \struck{\ill{}}
\marginal{cette \uncertain{\ill{}} \uncertain{marche}, plus géné., si $X\to S$ propre et lisse}
\uncertain{extension universelle} du schéma abélien dual $A$, et $E$ la
catégorie \uncertain{vectorielle} universelle $E$, la question est
équivalente à la suivante, \struck{\ill{}} :

c) Soit $\varphi$ une section de $E$ sur $S$, supposons que pour tout pt
\emph{fermé} $s\in S$, $\varphi(s)\in A_s^{(p(k(s)))}$ ($= pE_s$),
\uncertain{\ill{}} $p = $ car $k(s)$. Alors \struck{\ill{}} la \add{section}
$\varphi$ de $\Gamma(S,A_s)$ est-elle d'\emph{ordre fini} ? [\emph{NB} a)
signifie que l'image de $\varphi$ dans $\Gamma(S,A_s)$ est d'ordre fini, b)
signifie que si \add{\uncertain{\ill{} de plus que}} $\varphi$ est une section
de $V = \operatorname{Ker}(E\to A)$, alors $\varphi$ est nulle].

Ceci suggère la question

d) Soit $S$ de type fini sur $\mathbb{Z}$, \struck{\ill{}} \add{intègre,}
dominant $\operatorname{Spec}\mathbb{Z}$, \struck{\ill{}} $E$ un schéma en
groupes \struck{\ill{}} de tf. sur $S$ (\uncertain{lisse} \struck{connexe}
\uncertain{\ill{}}), $\omega$ une section de $E$. Supposons que pour tout pt
fermé $s\in S$, $\omega(s)$ appartienne \struck{\ill{}} à un sous-groupe
\add{connexe} de $E_s$ qui est \uncertain{\ill{}} \add{connexe} sur $k(s)$
[resp. \add{connexe} de rang unipotent nul, resp. \add{connexe} affine, resp.
\add{connexe} unipotent] alors $\omega(\eta)$ appartient à un sous-groupe de
$E_\eta$ qui est \uncertain{\ill{}}
\note{La phrase se poursuit p. 195, qui commence par « (resp. de rang unipotent nul, resp. affine, resp. unipotent) » : l'ordre des feuillets dans le dossier n'est pas celui de la rédaction.}

\page{197}
La partie a) dans c) suggère la question suivante :

e) Soit $A$ schéma abélien sur $S$, $\omega$ section de $A$, telle que pour
tout pt fermé $s\in S$, $\omega(s)$ se remonte à $A_s^{(p)}$ via
Verschiebung $A_s^{(p/S)}\to A_s$. Est-il vrai qu'elle est d'ordre fini ?

\struck{Dans le cas où $S$ \ill{} est de car.~$p>0$ (au lieu de dominer
$\operatorname{Spec}\mathbb{Z}$) on dispose de $A^{p/S}\to A$, et l'image
\ill{} de $\omega$ dans $A^{p/S}$ est un \uncertain{\ill{}}}

Fin 1ère partie.
\note{Sous ces lignes, la feuille est blanche ; on y voit par transparence l'écriture de la page suivante.}

\page{198}
\marginal{On suppose donc $f_*(\Omega^1_{X/S}) = f_*(Z^1_{X/S})$ et $\Lambda\simeq\mathcal{O}_{X^{p/S}}$}
\marginal{\emph{NB} au lieu de $f_*$, on pourrait prendre $R^if_*$}
Soit $X$ schéma sur $S$ tel que l'on soit sous les conditions permettant de
définir $\widetilde{C} : f_{X/S*}(\Omega^1_{X/S})\to\Omega^1_{X^{p/S}/S}$,
d'où
\[
  f_*(\Omega^1_{X/S}) = \underline{\omega} \xrightarrow{\ \mathrm{dfn}\ }
  g_*(\Omega^1_{X^{p/S}/S}) \simeq f_*(\Omega^1_{X/S})^{(p)} =
  \underline{\omega}^{(p)}.
\]
\note{Sous le $\simeq$, de sa main : « sous réserve que … ».}
Donc on trouve une \uncertain{flèche} $\mathcal{O}_S$-linéaire
\[
  \widetilde{C} : \underline{\omega}\longrightarrow\underline{\omega}^{(p)}
\]
on \uncertain{peut} former donc
\[
  \widetilde{C}^{(p^i)} : \underline{\omega}^{(p^i)} \longrightarrow
  (\underline{\omega}^{(p)})^{(p^i)} = \underline{\omega}^{(p^{i+1})},
\]
et on trouve un composé $\mathcal{O}_S$-linéaire
\[
  F_*^{(\nu)} = \underbrace{\widetilde{C}^{(p^{\nu-1})}\,
  \widetilde{C}^{(p^{\nu-2})}\cdots\widetilde{C}^{p}\,\widetilde{C}}_{\nu\ \text{facteurs}}
  : \underline{\omega}\longrightarrow\underline{\omega}^{(p^\nu)}
\]
Si $\lambda\mapsto\lambda^{p^\nu}$ sur $\mathcal{O}_S^*$ est l'identité,
\uncertain{par exemple} si $S$ est \uncertain{au-dessus} de $\mathbb{F}_{p^\nu}$,
on trouve donc que
\[
  F_*^{(\nu)} : \underline{\omega}\longrightarrow\underline{\omega}
\]
est un \emph{endomorphisme} \add{$\mathcal{O}_S$-linéaire} de
$\underline{\omega}$.

Soit $\omega$ une section de $\underline{\omega}$ telle que l'on ait
\[
  \widetilde{C}\omega = \omega^{(p)}
\]
(ce qui correspond à une \uncertain{section} $1$-forme $\omega$ sur $X$ telle
que $\widetilde{C}\omega = \widetilde{\omega}$ i.e. qui est loc. une
différentielle logarithmique …). On conclut par récurrence
\[
  \text{\struck{$\widetilde{C}^{(p^\nu)}$}}\ F^{(\nu)}\omega = \omega^{(p^\nu)}\,,
\]
et dans le cas où $\lambda = \lambda^{p^\nu}$ sur $\mathcal{O}_S^*$, on conclut
que l'on a
\[
  \boxed{\ F_*^{(\nu)}\omega = \omega\ }
\]

\page{199}
\struck{Soit $X$ schéma sur $S$, tel que $X\to X^{p/S}$ soit un
épimorphisme \uncertain{plat} (p.ex. $X$ lisse sur $S$, plus géné.
\ill{} $X\to S$ \ill{} \ill{} loc. \ill{} \ill{} \ill{} \ill{}, et
$X/S$, et les fibres géométriques \ill{}) et $\Omega^1_{X/S}$ est loc.
libre}
\note{Seul ce passage, encadré et barré, occupe la page.}

\page{200}
Notons que l'on peut aussi considérer
\[
  f_{X/S} : X\longrightarrow X^{p/S}
\]
d'où, par itération, des
\[
  f_{X/S}^{p^i/S} : X^{p^i/S}\longrightarrow X^{p^{i+1}/S}
\]
et par suite un composé
\[
  F_{X/S}^{(\nu)} = f_{X/S}^{(p^{\nu-1})}\,f_{X/S}^{(p^{\nu-2})}\cdots f_{X/S}^{p}\,f_{X/S}
  : X\longrightarrow X^{(p^\nu/S)}
\]
et on a, si $X/S$ est \emph{lisse de dim relative 1}
\[
  F_*^{(\nu)} = \operatorname{Tr}_{F^{(\nu)}_{X/S}} = \bigl(F^{\nu}_{X/S}\bigr)_*
\]
\note{Sous « Tr », le mot « Cartier » ; lecture incertaine.}
en identifiant (\uncertain{en vertu} du \uncertain{formalisme} de chgt. de
base) $g_{\nu*}(\Omega^1_{X^{p^\nu/S}})$ avec
$g_{0*}(\Omega^1_{X/S})^{(p^\nu)} = \underline{\omega}^{(p^\nu)}$. \uncertain{En}
particulier, si $\lambda\mapsto\lambda^{p^\nu}$ est l'identité sur
$\mathcal{O}_S^*$, donc $X^{p^\nu/S}\overset{\text{can}}{\simeq}X$, \uncertain{on
a}
\[
  F^{\nu}_{X/S} : X\longrightarrow X
\]
est un \uncertain{endomorphisme}, qui d'ailleurs n'est autre que l'itéré
\uncertain{\ill{}} de $F_X$ (\uncertain{remarquer} \uncertain{donc} que
$F_X^{\nu}$ est \uncertain{\ill{}} par hyp. $F_S^{\nu} = \mathrm{id}_S$, donc
\uncertain{que} $F_X^{\nu}$ est bien un $S$-endomorphisme) [En général, sans
hyp. sur $S$, on peut dire que le composé

\begin{tikzcd}
X \arrow[r, "F^{(\nu)}_{X/S}"] & X^{(p^\nu/S)} \arrow[r, "\text{can}"] \arrow[d, dashed] & X \arrow[d, dashed]\\
 & S \arrow[r, "F_S^{\nu}"'] & S
\end{tikzcd}

\note{Le carré de droite est marqué « cart » ; les flèches verticales sont pointillées sur la page.}

est $F_X^{\nu}$ ]
\note{L'argument se poursuit au-delà du lot.}

\end{document}
